\documentclass[11pt,a4paper]{article}

% --- Packages ---
\usepackage[utf8]{inputenc}
\usepackage[margin=1.0in]{geometry}
\usepackage{amsmath,amssymb}
\usepackage{graphicx}
\usepackage{booktabs}
\usepackage{multirow}
\usepackage{xcolor}
\usepackage{float}
\usepackage{subcaption}
\usepackage{hyperref}
\usepackage{fancyhdr}
\usepackage{array}

\hypersetup{
    colorlinks=true,
    linkcolor=blue,
    citecolor=blue,
    urlcolor=blue
}

% --- Header and Footer ---
\pagestyle{fancy}
\fancyhf{}
\rhead{\small Human Action Recognition for HRI}
\rfoot{\thepage}

\title{\Large \textbf{Human Action Recognition for Human-Robot Interaction (HRI)}}
\author{Alessandro Cortese \& Luca Rossetto}
\date{September 2026}

\begin{document}

\maketitle

\begin{abstract}
Reliable Human Action Recognition (HAR) is a fundamental prerequisite for safe and effective Human-Robot Interaction (HRI) in collaborative environments. This report presents a framework designed to detect and classify six fundamental human actions (\textit{boxing}, \textit{hand clapping}, \textit{hand waving}, \textit{jogging}, \textit{running}, and \textit{walking}) across continuous 40-frame image sequences from the KTH Action Recognition dataset. We implement and evaluate two complete pipelines: a \textbf{Classical Computer Vision} pipeline using KNN background subtraction, vertical morphological filtering, dense Farneb\"ack optical flow vector directionality, and SVM classification; and a \textbf{Deep Learning Fallback} pipeline using a pre-trained YOLOv8-Pose ONNX network for robust silhouette localization and pose-based feature extraction. Both pipelines feed into a multi-class Support Vector Machine (SVM) evaluated via 6-Fold Stratified CrossValidation. Experimental results demonstrate that while the classical approach achieves a Mean Intersection over Union (mIoU) of 0.7567 and a CV accuracy of 70.83\%, the deep learning fallback overcomes classical limits under lighting and clothing variations, achieving an mIoU of 0.8544 and a CV accuracy of 93.06\% (67 out of 72 sequences correctly classified).
\end{abstract}

\section{Introduction and Project Goal}
The main objective of this project is to develop a system capable of processing continuous 40-frame image sequences to perform two core tasks:
\begin{enumerate}
    \item \textbf{Subject Localization and Tracking }: Isolate the human actor from static background clutter and compute a dynamic bounding box for each frame, evaluated via Intersection over Union (IoU) on the 20th frame (median frame).
    \item \textbf{Kinematic Feature Extraction and Classification }: Analyze spatial-temporal motion dynamics across the 40 frame window, extract descriptive feature vectors, and classify the entire sequence with a single global semantic action label using a Support Vector Machine (SVM).
\end{enumerate}

The framework is tested against 72 stratified sequences from the benchmark KTH Action Recognition Dataset across six action classes: \textit{Boxing} (1), \textit{Handclapping} (2), \textit{Handwaving} (3), \textit{Jogging} (4), \textit{Running} (5), and \textit{Walking} (6).

\section{Dataset Stratification and Evaluation Protocol}
To validate system generalization without environmental bias, a stratified dataset of 72 video clips was curated. Each action class contains 12 sequences symmetrically sampled across four distinct KTH recording scenarios:
\begin{itemize}
    \item \textbf{Scenario $d1$}: Standard outdoor environment with uniform background and clear lighting.
    \item \textbf{Scenario $d2$}: Outdoor environment with scale variations, camera zoom shifts, and distance changes.
    \item \textbf{Scenario $d3$}: Outdoor environment with clothing variations.
    \item \textbf{Scenario $d4$}: Indoor environment with complex indoor illumination, low contrast, and cast shadows.
\end{itemize}

Each sequence folder contains exactly 40 pre extracted continuous image frames ($160 \times 120$ pixels) and a ground-truth annotation text file containing the numerical class label and the reference ground-truth bounding box coordinates for the median frame (frame 20):
\begin{equation}
\text{\texttt{<class\_id> <x\_center> <y\_center> <width> <height>}}
\end{equation}

\section{System Architecture and Software Design}
The code is written C++ using OpenCV and CMake. It follows an object-oriented design:

\begin{figure}[H]
    \centering
    \includegraphics[width=0.92\textwidth]{output/cv/person01_walking_d1.png}
    \caption{Example output visualization generated by the system, displaying the Ground Truth bounding box (blue), Predicted bounding box (green), predicted action label, and median frame IoU score.}
    \label{fig:example_output}
\end{figure}

\begin{itemize}
    \item \textbf{\texttt{DatasetLoader}}: Reads sequence frames and parses normalized groundtruth bounding box text files.
    \item \textbf{\texttt{Tracker} (CV Mode)}: Computes a sequence wide median background image $\mathbf{B}(x,y) = \text{median}_{t=1..N} \{\mathbf{I}_t(x,y)\}$ to eliminate moving subjects, applies absolute frame difference $|\mathbf{I}_t - \mathbf{B}|$, Gaussian smoothing, adaptive thresholding ($T=18 \to 10$), border clearing, and a $11 \times 25$ vertical morphological closing kernel to connect head, torso, and legs into a unified bounding box.
    \item \textbf{\texttt{YoloTracker} (DL Mode)}: Loads a lightweight YOLOv8-Pose ONNX network (\texttt{yolov8n-pose.onnx}) via OpenCV DNN module (\texttt{cv::dnn}). Detects 17 human body keypoints, computes their convex hull, expands the box to enclose wrists/ankles, and applies an Exponential Moving Average (EMA) filter ($\alpha=0.80$) across consecutive frames to suppress bounding box jitter.
    \item \textbf{\texttt{FeatureExtractor}} and \textbf{\texttt{YoloFeatureExtractor}}: Compute dense Farneb\"ack optical flow ($\mathbf{V}_x, \mathbf{V}_y$), body subregion flow energies (upper body vs lower body legs), percentile translation speeds, and centroid trajectory spans to produce a 27 element feature descriptor vector $\mathbf{x} \in \mathbb{R}^{27}$.
    \item \textbf{\texttt{ActionClassifier}}: Implements a multiclass Support Vector Machine (SVM).
    Hyperparameter tuning ($C \in [0.1, 500], \gamma \in [0.0001, 2.0]$) is performed using nested grid search, evaluated via 6-Fold Cross-Validation across all 72 sequences.
    \item \textbf{\texttt{main.cpp}}: Command-line entry point supporting three operational execution modes:
    \begin{enumerate}
        \item \texttt{./main} -- Runs Classical Computer Vision pipeline.
        \item \texttt{./main --use-yolo} -- Runs Deep Learning YOLOv8-Pose fallback pipeline.
        \item \texttt{./main --use-all} -- Runs both pipelines sequentially and prints a complete comparative summary table.
    \end{enumerate}
\end{itemize}

\section{Technical Challenges and Methodological Solutions}
Motion benchmark datasets present intrinsic physical ambiguities. Our feature extraction algorithms methodically address three primary technical challenges:

\subsection{Challenge I: Motion Directionality (Handclapping vs. Handwaving)}
\textit{Ambiguity}: Handclapping and Handwaving share near identical static human silhouettes, similar bounding box aspect ratios ($\text{width}/\text{height} \approx 0.40$), and virtually zero global centroid displacement ($X_{\text{span}} < 0.30$).

\textit{Solution}: We analyze local motion vector directionality within the upper body half ($Y \in [0, H/2]$):
\begin{enumerate}
    \item \textbf{Horizontal Vector Convergence (Clapping)}: We divide the upper body ROI into left and right halves. Clapping generates opposed horizontal optical flow vectors ($\mathbf{V}_x^{\text{left}} > 0$ and $\mathbf{V}_x^{\text{right}} < 0$). We compute average convergence speed:
    \begin{equation}
        \text{Conv}(t) = \frac{\bar{V}_{x,\text{left}}(t) - \bar{V}_{x,\text{right}}(t)}{H_{\text{actor}}}
    \end{equation}
    \item \textbf{Head/Hand Elevation Energy Ratio (Waving)}: Waving involves vertical and lateral arm oscillations above shoulder level. We extract optical flow energy in the top 30\% head region relative to total body energy ($E_{\text{head}} / E_{\text{total}}$), alongside vertical flow zero crossings.
\end{enumerate}

\subsection{Challenge II: Kinematic Intensity (Jogging vs. Running vs. Walking)}
\textit{Ambiguity}: Jogging, Running, and Walking exhibit identical leg stride mechanics and horizontal translation. Calculating total net displacement across the sequence ($\Delta X_{\text{net}} / \Delta t$) yields deceptive average values ($\approx 0.022$) because running actors traverse the scene in 12--18 frames while walking actors remain in frame for 35--40 frames.

\textit{Solution}: We measure \textbf{instantaneous frame to frame translation velocity} ($\Delta X_{\text{frame}} = |X_t - X_{t-1}| / H_{\text{actor}}$) and aggregate them using upper percentiles ($P_{75}$ and $P_{90}$):
\begin{itemize}
    \item \textbf{Walking}: $P_{75} \le 0.10$ px/frame (low instantaneous translation).
    \item \textbf{Jogging}: $0.15 < P_{75} \le 0.40$ px/frame (moderate translation).
    \item \textbf{Running}: $P_{75} > 0.50$ px/frame (high translation peaks up to 2.5 px/frame).
\end{itemize}
This is combined with lower body leg flow energy standard deviation ($\sigma_{\text{leg}}$) and gait zero crossings to capture stride rhythm frequency.

\subsection{Challenge III: Static vs. Dynamic Action Separation \& Arm Jitter}
\textit{Ambiguity}: When an actor punches in \textit{Boxing}, extending the arm expands the bounding box width ($W$), shifting the box center ($X + W/2$) by 10--15 pixels in a single frame. This arm extension pulse can trick a classifier into misinterpreting Boxing as horizontal translation.

\textit{Solution}: We apply a 3 frame moving average filter on bounding box centroids ($\mathbf{C}_t = \frac{\mathbf{C}_{t-1} + \mathbf{C}_t + \mathbf{C}_{t+1}}{3}$) to smooth out single frame arm pulses. Furthermore, we compute the \textbf{maximum normalized horizontal centroid trajectory span}:
\begin{equation}
    X_{\text{span}} = \frac{\max_{t} (X_t) - \min_{t} (X_t)}{H_{\text{actor}}}
\end{equation}
For static actions (Boxing, Clapping, Waving), $X_{\text{span}} < 0.45$, whereas for dynamic actions (Jogging, Running, Walking), $X_{\text{span}} > 0.85$, eliminating cross category static vs dynamic misclassifications.

\section{Limits of Classical CV vs. Deep Learning Advantages}
Evaluating both pipelines highlights fundamental trade offs between traditional algorithmic vision and modern deep learning:

\subsection{Limits of Classical Computer Vision}
\begin{enumerate}
    \item \textbf{Sensitivity to Contrast and Illumination ($d4$)}: In indoor scenes with shifting shadows, pixel intensity difference $|\mathbf{I}_t - \mathbf{B}|$ drops below fixed thresholds ($T=18$), causing foreground fragmentation and bounding box tracking loss.
    \item \textbf{Silhouette Tearing with Dark Clothing ($d3$)}: Dark coats blend into dark outdoor backgrounds, causing contour detection to return partial body boxes (legs or torso only), lowering median frame mIoU (0.7567).
    \item \textbf{Dependency on Background Staticity}: Background subtraction breaks down if tree leaves move or if minor camera vibrations occur.
\end{enumerate}

\subsection{Advantages of Pre-trained Deep Learning (YOLOv8-Pose)}
\begin{enumerate}
    \item \textbf{Invariant Human Pose Keypoint Localization}: YOLOv8-Pose detects 17 keypoints (nose, shoulders, elbows, wrists, hips, knees, ankles) based on learned spatial feature pyramids, completely invariant to background lighting ($d4$) or clothing ($d3$).
    \item \textbf{Convex Hull Mask Generation}: Connecting detected keypoints into a convex polygon produces clean human masks even when background subtraction fails.
    \item \textbf{Higher Localization Precision}: Increases Mean IoU from 0.7567 to \textbf{0.8544} and 6-Fold Cross-Validation accuracy from 70.83\% to \textbf{93.06\%}.
\end{enumerate}

\section{Experimental Results and Comparative Analysis}

\subsection{Localization Accuracy (mIoU)}
Bounding box localization is evaluated on the 20th frame (median frame) against annotated ground-truth boxes using Intersection over Union:
\begin{equation}
    \text{IoU} = \frac{\text{Area}(\text{Box}_{\text{pred}} \cap \text{Box}_{\text{gt}})}{\text{Area}(\text{Box}_{\text{pred}} \cup \text{Box}_{\text{gt}})}
\end{equation}

\begin{table}[H]
\centering
\caption{Per class and Mean IoU (mIoU) comparison between Classical CV and YOLOv8-Pose.}
\label{tab:miou_comparison}
\begin{tabular}{lcc}
\toprule
\textbf{Action Class} & \textbf{Classical CV (mIoU)} & \textbf{Deep Learning YOLOv8 (mIoU)} \\
\midrule
Boxing & 0.6676 & 0.8050 \\
Clapping & 0.7155 & 0.8938 \\
Waving & 0.5780 & 0.9121 \\
Jogging & 0.7627 & 0.8176 \\
Running & 0.9062 & 0.8596 \\
Walking & 0.9095 & 0.8381 \\
\midrule
\textbf{Overall mIoU} & \textbf{0.7566} & \textbf{0.8544} \\
\bottomrule
\end{tabular}
\end{table}

\subsection{Classification Accuracy and F1-Scores}
Classification performance is evaluated using Stratified 6-Fold Cross-Validation across all 72 sequences.

\begin{table}[H]
\centering
\small
\caption{Stratified 6Fold Cross-Validation classification accuracy and perclass F1 scores.}
\label{tab:classification_comparison}
\begin{tabular}{l@{\hspace{8pt}}c@{\hspace{8pt}}c@{\hspace{8pt}}c@{\hspace{8pt}}c@{\hspace{8pt}}c@{\hspace{8pt}}c}
\toprule
& \multicolumn{3}{c}{\textbf{Classical CV (Acc: 70.83\%)}} & \multicolumn{3}{c}{\textbf{Deep Learning YOLO (Acc: 93.06\%)}} \\
\cmidrule(r){2-4} \cmidrule(l){5-7}
\textbf{Action Class} & \textbf{Precision} & \textbf{Recall} & \textbf{F1-Score} & \textbf{Precision} & \textbf{Recall} & \textbf{F1-Score} \\
\midrule
Boxing & 0.56 & 0.75 & 0.64 & 1.00 & 0.92 & \textbf{0.96} \\
Clapping & 0.71 & 0.83 & 0.77 & 0.92 & 1.00 & \textbf{0.96} \\
Waving & 0.70 & 0.58 & 0.64 & 0.92 & 0.92 & \textbf{0.92} \\
Jogging & 0.62 & 0.67 & 0.64 & 0.80 & 1.00 & \textbf{0.89} \\
Running & 0.78 & 0.58 & 0.67 & 1.00 & 0.75 & \textbf{0.86} \\
Walking & 1.00 & 0.83 & 0.91 & 1.00 & 1.00 & \textbf{1.00} \\
\midrule
\textbf{Overall Mean} & \textbf{0.73} & \textbf{0.71} & \textbf{0.71} & \textbf{0.94} & \textbf{0.93} & \textbf{0.93} \\
\bottomrule
\end{tabular}
\end{table}

\subsection{Confusion Matrices}
Comparing the cumulative confusion matrices illustrates how YOLO pose tracking eliminates cross-category misclassifications:

\begin{figure}[H]
\centering
\begin{subfigure}[b]{0.48\textwidth}
\centering
\begin{verbatim}
      [1]  [2]  [3]  [4]  [5]  [6]
[1]     9    2    1    0    0    0
[2]     1   10    1    0    0    0
[3]     3    2    7    0    0    0
[4]     1    0    1    8    2    0
[5]     1    0    0    4    7    0
[6]     1    0    0    1    0   10
\end{verbatim}
\caption{Classical CV Matrix (Acc: 70.83\%)}
\end{subfigure}
\hfill
\begin{subfigure}[b]{0.48\textwidth}
\centering
\begin{verbatim}
      [1]  [2]  [3]  [4]  [5]  [6]
[1]    11    0    1    0    0    0
[2]     0   12    0    0    0    0
[3]     0    1   11    0    0    0
[4]     0    0    0   12    0    0
[5]     0    0    0    3    9    0
[6]     0    0    0    0    0   12
\end{verbatim}
\caption{Deep Learning YOLO Matrix (Acc: 93.06\%)}
\end{subfigure}
\caption{Cumulative 6-Fold Confusion Matrices (1: Boxing, 2: Clapping, 3: Waving, 4: Jogging, 5: Running, 6: Walking).}
\label{fig:confusion_matrices}
\end{figure}

In the Deep Learning pipeline, 67 out of 72 sequences (93.06\%) are perfectly classified. The only remaining minor confusion occurs between Jogging (4) and Running (5), where 3 running sequences are classified as jogging due to speed overlaps in actor $p02$.

\section{Sample Output Visualizations}
Figure~\ref{fig:all_classes_viz} shows visual outputs generated across the six action classes under different environmental scenarios ($d1, d2, d3, d4$). In each image, the Ground Truth bounding box is drawn in blue, the predicted bounding box in green, with the predicted category label and median frame IoU displayed.

\begin{figure}[H]
    \centering
    \begin{subfigure}[b]{0.32\textwidth}
        \centering
        \includegraphics[width=\textwidth]{output/yolo/person01_boxing_d1.png}
        \caption{Boxing ($d1$)}
    \end{subfigure}
    \hfill
    \begin{subfigure}[b]{0.32\textwidth}
        \centering
        \includegraphics[width=\textwidth]{output/yolo/person01_handclapping_d1.png}
        \caption{Handclapping ($d1$)}
    \end{subfigure}
    \hfill
    \begin{subfigure}[b]{0.32\textwidth}
        \centering
        \includegraphics[width=\textwidth]{output/yolo/person01_handwaving_d1.png}
        \caption{Handwaving ($d1$)}
    \end{subfigure}

    \vspace{0.5em}

    \begin{subfigure}[b]{0.32\textwidth}
        \centering
        \includegraphics[width=\textwidth]{output/yolo/person01_jogging_d1.png}
        \caption{Jogging ($d1$)}
    \end{subfigure}
    \hfill
    \begin{subfigure}[b]{0.32\textwidth}
        \centering
        \includegraphics[width=\textwidth]{output/yolo/person01_running_d1.png}
        \caption{Running ($d1$)}
    \end{subfigure}
    \hfill
    \begin{subfigure}[b]{0.32\textwidth}
        \centering
        \includegraphics[width=\textwidth]{output/yolo/person01_walking_d1.png}
        \caption{Walking ($d1$)}
    \end{subfigure}
    \caption{Sample target outputs produced by the Deep Learning YOLO pipeline for all six action categories.}
    \label{fig:all_classes_viz}
\end{figure}

\section{Team Member Contributions and Working Hours}

\subsection{Rossetto Luca (16 Hours)}
\begin{itemize}
    \item \textbf{Ideas \& System Architecture}: Codesigned the dual pipeline system framework comparing Classical Computer Vision against Deep Learning fallback tracking.
    \item \textbf{Implementation}: Developed the Data Loading and Bounding Box Localization modules (\texttt{DatasetLoader}, \texttt{Tracker}, and \texttt{YoloTracker}), including median background subtraction, adaptive thresholding, morphological filtering, YOLOv8-Pose ONNX network integration, and bounding box drawing routines.
    \item \textbf{Testing}: Conducted component-level tests for frame sequence loading, background model stability, and bounding box localization across all environmental scenarios ($d1, d2, d3, d4$).
    \item \textbf{Performance Measurement}: Computed and analyzed tracking performance metrics, evaluating frame-level and median Intersection over Union (mIoU) scores.
\end{itemize}

\subsection{Cortese Alessandro (16 Hours)}
\begin{itemize}
    \item \textbf{Ideas \& System Architecture}: Formulated the 27-element spatiotemporal feature descriptor vector (optical flow directionality, energy distribution, relative scale normalization) and SVM classification scheme.
    \item \textbf{Implementation}: Developed Feature Extraction and Classification (\texttt{FeatureExtractor}, \texttt{YoloFeatureExtractor}, and \texttt{ActionClassifier}), leveraging Farneb\"ack optical flow calculation, dynamic velocity scaling, and RBF kernel Support Vector Machine.
    \item \textbf{Testing}: Built the Stratified 6-Fold Cross-Validation pipeline, Z-score feature standardization, and grid search hyperparameter tuning.
    \item \textbf{Performance Measurement}: Evaluated classification metrics across all 72 sequences, computing per class Precision, Recall, F1-Scores, overall accuracy, and confusion matrices.
\end{itemize}

\subsection{Joint Contributions}
Both team members jointly collaborated on end-to-end system integration in \texttt{main.cpp} (including the \texttt{--use-all} comparative execution mode), code refactoring, and authoring the final technical LaTeX report.

\end{document}
